\section{线路分段、易损度、故障抽样与修复}
\label{sec:fragility-fault-repair}

\subsection{自动分段与资产类型}
\srcpath{generate\_typhoon\_line\_segments} 遍历 AC/DC branches。若有母线坐标，沿端点线性插值
构造段；否则围绕 case center 按 bus ID 派生 fallback 坐标，并置
\fld{used\_fallback\_coordinates}。段数近似为支路长度除目标段长，再钳在
\fld{min\_segments\_per\_branch} 与 \fld{max\_segments\_per\_branch}。无法得到实际几何时置
\fld{used\_synthetic\_segments}。

段资产类型为 OverheadLine、Cable 或 Mixed。案例缺少元数据时按
\fld{default\_overhead\_fraction} / \fld{default\_cable\_fraction} 选择默认易损参数；这不是从
线路名称猜测结构类型。

\subsection{风雨等效风速}
架空段先将阵风和降雨组合。令 $v_m=V/1.42$：
\begin{align}
 f_1&=e^{0.006462v_m}-1.2486e^{-0.2769v_m},\\
 f_2&=\begin{cases}0.09376R^{0.7087},&R>0,\\0,&R=0,\end{cases}\\
 V_{eq}&=\max(0,(v_m+\max(0,f_1f_2))1.42).
\end{align}

\subsection{对数正态风易损度}
当 $V_{eq}\le V_d$ 时 $p_w=0$。否则
\begin{align}
 z&=\frac{\ln V_{eq}-\ln V_{50}}{\sigma_{\ln V}},\\
 s&=\begin{cases}
 1,&V_c>0\land V_{eq}\ge V_c,\\
 \Phi(z),&\text{其他},
 \end{cases}\\
 \ell_e&=\operatorname{clip}(\ell,0,5),\\
 p_w&=\operatorname{clip}\left(1-e^{-0.015s\ell_e\Delta t},0,1\right).
\end{align}
$V_c=$\fld{collapse\_wind\_ms} 的饱和分支防止超高风速下数值尾部行为逆转。非正的中位风速或
sigma 由底层 CDF 防御为零 severity。

\subsection{风雨联合段暴露}
令 $r_V=V/V_d$（$V_d\le0$ 时为 0）、$r_R=R/R_d$（$R_d\le0$ 时为 0）：
\begin{align}
 \lambda_s&=\ell^2\exp(a r_V+b r_R+c),\\
 p_s&=\begin{cases}
 \operatorname{clip}(1-e^{-\lambda_s\Delta t},0,1),&V>V_d\lor R>R_d,\\
 0,&\text{其他},
 \end{cases}\\
 p_o&=\operatorname{clip}(1-(1-p_w)(1-p_s),0,1).
\end{align}
该合成假设两个条件机制独立，仅是模型结构，不是由数据估计出的独立性结论。

\subsection{电缆水文与失效概率}
对 Curve Number $CN$：
\begin{align}
 CN'&=\operatorname{clip}(CN,1,99),\quad S=25400/CN'-254,\quad I_a=0.2S,\\
 Q&=\begin{cases}(R-I_a)^2/\max(10^{-9},R-I_a+S),&R>I_a,\\0,&R\le I_a,\end{cases}\\
 I&=\max(0,R-Q),\quad q=Q/(1000\cdot3600),\\
 h&=\left(\frac{q n}{\sqrt{\max(10^{-9},S_0)}}\right)^{0.6},\\
 \theta&=\min\left(\phi,\theta_0+\frac{I}{\max(10^{-9},1000Z)}\right),\\
 p_c&=\operatorname{clip}\left(\frac{p_{c,0}}{10}
 e^{130h+8\theta/\max(10^{-9},\phi)},0,1\right).
\end{align}
Cable 返回 $p_c$，Overhead 返回 $p_o$，Mixed 返回 $1-(1-p_o)(1-p_c)$ 后再钳制。

\subsection{段到支路聚合}
同一支路各段按条件独立并联失效：
\begin{equation}
 p_{b,t}=\operatorname{clip}\left(1-\prod_{s\in b}(1-p_{s,t}),0,1\right).
\end{equation}
结果逐支路保存峰值风速、峰值雨强和峰值失效概率。段越多会提高并联暴露，因此目标段长会
改变支路概率；必须与易损参数共同校准。

\subsection{抽样支持门}
对尚未故障的支路逐步生成 $U\sim U(0,1)$。只有
\begin{equation}
 p_{b,t}\ge p_{min}\quad\text{且}\quad U<p_{b,t}
\end{equation}
时记录首次故障。\fld{min\_fault\_probability} 是抽样支持门，不是风险报告过滤器：低于门槛
的 $p_{b,t}$ 仍进入 branch risk。固定案例把 gate 设为 1，实测峰值概率 0.798910，因此
fault count 为 0，但风险保留。

\subsection{修复时长}
基础 MTTR 优先取资产的正值，否则用 \fld{default\_repair\_time\_hr}：
\begin{align}
 T_r'&=T_r\left[1+k_w\frac{\max(0,V_{peak}-25)}{25}\right],\\
 T_r''&=\operatorname{clip}(T_r',\max(10^{-6},\Delta t),T_{r,max}).
\end{align}
故障事件保存 fault hour 和 repair hour，供下游恢复模块消费。

\subsection{分阶段抢修}
准备时刻为最后故障时刻加非负 delay；维修槽宽
\begin{equation}
 \Delta t_c=\max(\Delta t,\texttt{repair\_crew\_time\_per\_branch\_hr}).
\end{equation}
在精确预算内，每轮试验恢复一条故障支路，依据恢复供电增量及确定性破同分项：
\begin{equation}
 score=impact+10^{-3}p_{peak}+10^{-5}V_{peak}+10^{-6}\ell.
\end{equation}
这是一阶段贪心，不是全时域组合最优。当故障数超过内部精确预算，改按源距离、故障时间、
峰值概率、支路类型和稳定 ID 排序，置 \fld{used\_approximate\_repair\_order=true}。

\subsection{算法流程}
\begin{enumerate}
  \item 生成/接收线路分段；对每个轨迹步计算段中点风雨；
  \item 依资产类型计算段概率并聚合为支路概率；
  \item 更新支路峰值风险，应用 sampling gate 和首次故障抽样；
  \item 计算 MTTR 风速修正；需要时执行 staged repair；
  \item 生成 renewable profile multiplier 和完整 validity flags。
\end{enumerate}

\subsection{复杂度和边界}
风雨/易损链为 $O(TL_s)$。精确 staged repair 的供电增量试验可达 $O(F^2G)$，其中 $F$ 为故障
数、$G$ 为图评估成本；超过预算后排序为 $O(F\log F)$。

\begin{gapnote}
当前没有杆塔级相关失效、飞散物路径、保护动作、维修队空间路由、多队并行资源约束或全局
最优修复调度。分段独立性、默认易损参数和 sampling gate 都必须作为案例假设披露。
\end{gapnote}

\subsection{实现锚点}
\begin{center}\small
\begin{tabular}{@{}L{5.1cm}L{8.6cm}@{}}
\toprule
对象 & 源码锚点\\
\midrule
自动分段 & \srcpath{typhoon\_resilience.cpp:generate\_typhoon\_line\_segments}\\
架空/电缆概率 & \srcpath{overhead\_segment\_probability}、
\srcpath{cable\_segment\_probability}\\
支路聚合与抽样 & \srcpath{generate\_typhoon\_fault\_sequence}\\
修复排序 & \srcpath{apply\_staged\_post\_disaster\_repairs}\\
\bottomrule
\end{tabular}
\end{center}
